% Generated draft: requires mathematical and source review.
\documentclass[11pt]{article}
\usepackage[margin=25mm]{geometry}
\usepackage{fontspec}
\setmainfont{Latin Modern Roman}
\usepackage{amsmath,amssymb,mathtools,unicode-math}
\setmathfont{Latin Modern Math}
\usepackage{enumitem,needspace,xurl,hyperref}
\setlength{\parindent}{0pt}
\setlength{\parskip}{5pt}
\setlength{\emergencystretch}{3em}
\newcommand{\N}{\mathbb{N}}
\newcommand{\Z}{\mathbb{Z}}
\newcommand{\Q}{\mathbb{Q}}
\newcommand{\R}{\mathbb{R}}
\newcommand{\C}{\mathbb{C}}
\newcommand{\F}{\mathbb{F}}
\newcommand{\A}{\mathbb{A}}
\newcommand{\PP}{\mathbb{P}}
\newcommand{\E}{\mathbb{E}}
\newcommand{\eps}{\varepsilon}
\newcommand{\dd}{\,\mathrm{d}}
\begin{document}
\section*{UnsolvedMath ID: 30001304}
\textbf{Source problem code:} catalogue entry
\subsection{Definitions and mathematical statement}
Under the usual meaning of \emph{almost totally real} (ATR), let \(F\) be a totally real number field of degree \(d\), and let \(K/F\) be a quadratic extension such that exactly one real place of \(F\) becomes complex in \(K\), while every other real place splits into two real places. Equivalently, \(K\) has signature \((2d-2,1)\). The supplied record does not specify any further degree restriction or arithmetic hypotheses.

The source considers an \(F\)-algebra embedding of \(K\) into \(M_2(F)\), the algebra of \(2\times2\) matrices over \(F\). Retain its notation \(E\), \(f\), \(\Psi\), \(\Delta_\Psi\), \(AJ_f\), \(c_\Psi\), and \(H_{c_\Psi}\). Write
\[
P_\Psi=AJ_f(\Delta_\Psi)\in E(\C)
\]
for the associated complex ATR point whenever the source construction is defined, using the complex realization implicit in that construction.

\textbf{Conjecture as supplied.} For every \(\Psi\) admitted by the source construction, after choosing the isogeny in the definition of \(AJ_f\) as required there, the associated point satisfies
\[
AJ_f(\Delta_\Psi)\in E(H_{c_\Psi}).
\]
In particular, the source asserts that this complex ATR point is algebraic.

\textbf{Unresolved qualifications.} The record does not define the elliptic curve \(E\), the form \(f\), the admissible embeddings or other objects denoted by \(\Psi\), the cycle \(\Delta_\Psi\), the map \(AJ_f\), the parameter \(c_\Psi\), or the field \(H_{c_\Psi}\). It also does not specify the required isogeny, whether a single choice must work for every \(\Psi\), or the field embeddings used to interpret the displayed inclusion. In particular, the meaning and normalization of the field notation cannot be fixed from this record. Thus the displayed assertion preserves the supplied conjecture but remains a statement schema pending recovery of these definitions and hypotheses from the cited source.
\subsection{Short English statement}
The supplied conjecture says that every complex ATR point arising from the intended construction is rational over the designated field \(H_{c_\Psi}\), after the isogeny choice required in defining \(AJ_f\). The construction, its admissible inputs, and the field are insufficiently specified in the record for a fully standalone formulation.
\subsection{Sources}
\begin{itemize}
\item Local UnsolvedMath registry, supplied complete source record for permanent ID \(30001304\), title \emph{Algebraicity of Complex ATR Points}, exact source problem code \texttt{catalogue entry}, in the Oberwolfach Reports Open Problems collection. The record gives proposed year 2009 and supplies no named proposer.
\item Supplied report citation: \emph{Explicit Methods in Number Theory (2010)}, Oberwolfach Reports, DOI \texttt{10.4171/owr/2009/33}, pp.~1843--1920. \url{https://doi.org/10.4171/owr/2009/33}. The record describes Oberwolfach Reports as documenting mathematical workshops and their open-problem discussions. The year 2010 in the citation and the 2009 provenance in the record are retained without reconciliation.
\item Supplied literature source, labeled ``Darmon Oberwolfach report'': \url{https://ems.press/content/serial-article-files/46233?nt=1}.
\item Supplied literature source, labeled ``ATR Darmon points computation'': \url{https://arxiv.org/abs/1204.6680}.
\item Supplied literature source, labeled ``Stark--Heegner points'': \url{https://arxiv.org/abs/2207.01310}.
\item Supplied literature assessment dated exactly \texttt{2026-08-21}: the registry labels the problem \texttt{partially\_solved} and states that ATR Darmon-point algebraicity remains conjectural in general, with numerical computations and special geometrically modular cases providing evidence and partial results. This is supplied dated curation derived from the OWR material, not an independently verified literature assessment; its time-sensitive status and cited claims require independent verification.
\item The supplied statement status is \texttt{corrected\_verified}; its verification note describes meaning-preserving semantic, typographic, and LaTeX cleanup with no unsupported source-dependent additions. These are registry annotations, not independent verification of this draft or of the conjecture.
\end{itemize}
\end{document}
